\chapter{2013年新课标I卷（文）}

\section{单选题}

\begin{problem}
已知集合 \(A = \{1, 2, 3, 4\}\)，\(B = \{x \mid x = n^2, n \in A\}\)，则 \(A \cap B =\)（\quad）

\choices
    {\(\{1, 4\}\)}
    {\(\{2, 3\}\)}
    {\(\{9, 16\}\)}
    {\(\{1, 2\}\)}
\end{problem}

\begin{answer}
A
\end{answer}

\begin{solution}
由 \(n \in A\) 可得 \(B=\{1^2,2^2,3^2,4^2\}=\{1,4,9,16\}\)，因此 \(A \cap B=\{1,4\}\)，故选 A.
\end{solution}

\begin{problem}
\(\frac{1 + 2\i}{(1 - \i)^2} =\)（\quad）

\choices
    {\(-1 - \frac{1}{2}\i\)}
    {\(-1 + \frac{1}{2}\i\)}
    {\(1 + \frac{1}{2}\i\)}
    {\(1 - \frac{1}{2}\i\)}
\end{problem}

\begin{answer}
B
\end{answer}

\begin{solution}
先计算分母：\((1-\i)^2=1-2\i+\i^2=-2\i\). 因而
\(\displaystyle\frac{1+2\i}{(1-\i)^2}=\frac{1+2\i}{-2\i}=-1+\frac{1}{2}\i\)，故选 B.
\end{solution}

\begin{problem}
从 \(1\)，\(2\)，\(3\)，\(4\) 中任取 \(2\) 个不同的数，则取出的 \(2\) 个数之差的绝对值为 \(2\) 的概率是（\quad）

\choices
    {\(\frac{1}{2}\)}
    {\(\frac{1}{3}\)}
    {\( \frac{1}{4} \)}
    {\(\frac{1}{6}\)}
\end{problem}

\begin{answer}
B
\end{answer}

\begin{solution}
从 \(1\)，\(2\)，\(3\)，\(4\) 中任取两个不同的数共有 \(\mathrm C_4^2=6\) 种取法. 差的绝对值为 \(2\) 的取法只有 \(\{1,3\}\) 和 \(\{2,4\}\)，共有 \(2\) 种，所以所求概率为 \(\displaystyle\frac{2}{6}=\frac{1}{3}\)，故选 B.
\end{solution}

\begin{problem}
已知双曲线 \(C: \frac{x^2}{a^2} - \frac{y^2}{b^2} = 1\left(a > 0, b > 0\right)\) 的离心率为 \(\frac{\sqrt{5}}{2}\)，则 \(C\) 的渐近线方程为（\quad）

\choices
    {\(y = \pm \frac{1}{4} x\)}
    {\(y = \pm \frac{1}{3} x\)}
    {\(y = \pm \frac{1}{2} x\)}
    {\(y = \pm x\)}
\end{problem}

\begin{answer}
C
\end{answer}

\begin{solution}
双曲线的焦半距满足 \(c^2=a^2+b^2\)，且离心率 \(e=\frac{c}{a}=\frac{\sqrt{5}}{2}\). 因此
\(\displaystyle\frac{b^2}{a^2}=\frac{c^2-a^2}{a^2}=e^2-1=\frac{5}{4}-1=\frac{1}{4}\)，从而 \(\frac{b}{a}=\frac{1}{2}\). 双曲线的渐近线为 \(y=\pm\frac{b}{a}x\)，即 \(y=\pm\frac{1}{2}x\)，故选 C.
\end{solution}

\begin{problem}
已知命题 \(p\)：\(\forall x \in \R\)，\(2^x < 3^x\)；命题 \(q\)：\(\exists x \in \R\)，\(x^3 = 1 - x^2\)，则下列命题中为真命题的是（\quad）

\choices
    {\(p \wedge q\)}
    {\(\neg p \land q\)}
    {\(p \wedge \neg q\)}
    {\(\neg p \land \neg q\)}
\end{problem}

\begin{answer}
B
\end{answer}

\begin{solution}
命题 \(p\) 不成立，因为取 \(x=-1\) 时 \(2^{-1}=\frac{1}{2}>\frac{1}{3}=3^{-1}\). 令 \(h(x)=x^3+x^2-1\)，则 \(h(0)=-1<0\)，\(h(1)=1>0\)，由零点存在性定理，方程 \(x^3=1-x^2\) 在 \((0,1)\) 内有实根，所以命题 \(q\) 成立. 因此真命题为 \(\neg p\wedge q\)，故选 B.
\end{solution}

\begin{problem}
设首项为 \(1\)，公比为 \(\frac{2}{3}\) 的等比数列 \(\{a_n\}\) 的前 \(n\) 项和为 \(S_n\)，则（\quad）

\choices
    {\(S_{n} = 2a_{n} - 1\)}
    {\(S_{n} = 3a_{n} - 2\)}
    {\(S_{n} = 4 - 3a_{n}\)}
    {\(S_{n} = 3 - 2a_{n}\)}
\end{problem}

\begin{answer}
D
\end{answer}

\begin{solution}
等比数列的通项为 \(a_n=\left(\frac{2}{3}\right)^{n-1}\)，所以
\(\displaystyle S_n=\frac{1-\left(\frac{2}{3}\right)^n}{1-\frac{2}{3}}=3-3\left(\frac{2}{3}\right)^n=3-2\left(\frac{2}{3}\right)^{n-1}=3-2a_n\)，故选 D.
\end{solution}

\begin{problem}
执行下面的程序框图，若输入的 \(t \in [-1,3]\)，则输出的 \(s\) 属于（\quad）

\bitmapfigure[max width=0.26\linewidth]{img/2013/new_curriculum_paper_1_liberal/q7_fig1.png}

\choices
    {\([-3, 4]\)}
    {\([-5, 2]\)}
    {\([-4,3]\)}
    {\([-2,5]\)}
\end{problem}

\begin{answer}
A
\end{answer}

\begin{solution}
当 \(t<1\) 时，\(t\in[-1,1)\)，程序输出 \(s=3t\)，故 \(s\in[-3,3)\). 当 \(t\geqslant1\) 时，\(t\in[1,3]\)，输出 \(s=4t-t^2=4-(t-2)^2\)，故 \(s\in[3,4]\). 合并得 \(s\in[-3,4]\)，故选 A.
\end{solution}

\begin{problem}
\(O\) 为坐标原点，\(F\) 为抛物线 \(C: y^{2} = 4\sqrt{2} x\) 的焦点，\(P\) 为 \(C\) 上一点，若 \(|PF| = 4\sqrt{2}\)，则 \(\triangle POF\) 的面积为（\quad）

\choices
    {\(2\)}
    {\(2\sqrt{2}\)}
    {\(2\sqrt{3}\)}
    {\(4\)}
\end{problem}

\begin{answer}
C
\end{answer}

\begin{solution}
抛物线方程为 \(y^2=4px\)，所以 \(p=\sqrt{2}\)，焦点为 \(F(\sqrt{2},0)\). 设 \(P(x,y)\) 在抛物线上，由抛物线的焦半径公式得 \(PF=x+p\)，于是 \(x+\sqrt{2}=4\sqrt{2}\)，即 \(x=3\sqrt{2}\). 再由抛物线方程得 \(y^2=4\sqrt{2}\cdot3\sqrt{2}=24\)，所以 \(|y|=2\sqrt{6}\). 三角形 \(POF\) 以 \(OF=\sqrt{2}\) 为底，点 \(P\) 到 \(x\) 轴的距离为 \(2\sqrt{6}\)，故
\(\displaystyle S_{\triangle POF}=\frac{1}{2}\cdot\sqrt{2}\cdot2\sqrt{6}=2\sqrt{3}\)，故选 C.
\end{solution}

\begin{problem}
函数 \(f(x) = (1 - \cos x)\sin x\) 在 \(\left[-\pi,\pi\right]\) 的图象大致为（\quad）

\choices
    {\choicebitmap[width=0.15\paperwidth]{img/2013/new_curriculum_paper_1_liberal/q9_choice_a.png}}
    {\choicebitmap[width=0.15\paperwidth]{img/2013/new_curriculum_paper_1_liberal/q9_choice_b.png}}
    {\choicebitmap[width=0.15\paperwidth]{img/2013/new_curriculum_paper_1_liberal/q9_choice_c.png}}
    {\choicebitmap[width=0.15\paperwidth]{img/2013/new_curriculum_paper_1_liberal/q9_choice_d.png}}
\end{problem}

\begin{answer}
C
\end{answer}

\begin{solution}
在区间 \([-\pi,\pi]\) 上，\(1-\cos x\geqslant0\)，所以 \(f(x)\) 的符号与 \(\sin x\) 相同：在 \((-\pi,0)\) 上为负，在 \((0,\pi)\) 上为正，并且在 \(-\pi\)，\(0\)，\(\pi\) 处为零. 此外，\(f(-x)=-f(x)\)，图象关于原点对称. 又 \(f'(0)=0\)，且 \(f(x)\sim\frac{x^3}{2}\)（\(x\) 趋近于 \(0\)），所以图象在原点附近平缓地穿过 \(x\) 轴. 只有选项 C 同时符合这些特征，故选 C.
\end{solution}

\begin{problem}
已知锐角三角形 \(ABC\) 的内角 \(A\)，\(B\)，\(C\) 的对边分别为 \(a\)，\(b\)，\(c\)，\(23\cos^{2} A + \cos 2A = 0\)，\(a = 7\)，\(c = 6\)，则 \(b =\)（\quad）

\choices
    {\(10\)}
    {\(9\)}
    {\(8\)}
    {\(5\)}
\end{problem}

\begin{answer}
D
\end{answer}

\begin{solution}
由倍角公式 \(\cos 2A=2\cos^2A-1\)，原条件化为 \(23\cos^2A+2\cos^2A-1=0\)，即 \(25\cos^2A=1\). 因为三角形为锐角三角形，\(\cos A>0\)，所以 \(\cos A=\frac{1}{5}\). 由余弦定理，
\(\displaystyle a^2=b^2+c^2-2bc\cos A\)，代入 \(a=7\)，\(c=6\) 得 \(49=b^2+36-\frac{12}{5}b\)，即 \(5b^2-12b-65=0\). 解得 \(b=5\) 或 \(b=-\frac{13}{5}\). 因为边长为正，所以 \(b=5\)，故选 D.
\end{solution}

\begin{problem}
某几何体的三视图如图所示，则该几何体的体积为（\quad）

\bitmapfigure[max width=0.44\linewidth]{img/2013/new_curriculum_paper_1_liberal/q11_fig1.png}

\choices
    {\(16 + 8\pi\)}
    {\(8 + 8\pi\)}
    {\(16 + 16\pi\)}
    {\(8 + 16\pi\)}
\end{problem}

\begin{answer}
A
\end{answer}

\begin{solution}
由三视图可读出，几何体由一个长、宽、高分别为 \(4\)，\(2\)，\(2\) 的长方体和一个沿长为 \(4\) 的方向延伸、底面直径为 \(4\) 的半圆柱组成. 长方体体积为 \(4\times2\times2=16\)，半圆柱的底面半径为 \(2\)，体积为 \(\displaystyle\frac{1}{2}\pi\times2^2\times4=8\pi\). 所以几何体体积为 \(16+8\pi\)，故选 A.
\end{solution}

\begin{problem}
已知函数 \(f(x) = \begin{cases}
-x^2 + 2x, & x \leqslant 0,\\
\ln (x + 1), & x > 0,
\end{cases} \) 若 \(|f(x)| \geqslant ax\)，则 \( a \) 的取值范围是（\quad）

\choices
    {\(\left(-\infty,0\right]\)}
    {\(\left(-\infty,1\right]\)}
    {\(\left[-2,1\right]\)}
    {\(\left[-2,0\right]\)}
\end{problem}

\begin{answer}
D
\end{answer}

\begin{solution}
当 \(x\leqslant0\) 时，\(f(x)=-x^2+2x\leqslant0\)，所以 \(|f(x)|=x^2-2x\). 要使不等式对所有 \(x\leqslant0\) 成立，需要
\(x^2-2x\geqslant ax\)，即 \(x\bigl(x-a-2\bigr)\geqslant0\). 对全部 \(x\leqslant0\) 成立的充要条件是 \(a\geqslant-2\). 当 \(x>0\) 时，\(f(x)=\ln(x+1)>0\)，不等式等价于 \(\displaystyle a\leqslant\frac{\ln(x+1)}{x}\). 由于 \(\displaystyle\lim_{x\to+\infty}\frac{\ln(x+1)}{x}=0\)，要对所有 \(x>0\) 成立必须 \(a\leqslant0\)；反之 \(a\leqslant0\) 时右端 \(ax\leqslant0<f(x)\)，不等式确实成立. 因此 \(a\in[-2,0]\)，故选 D.
\end{solution}

\section{填空题}

\begin{problem}
已知两个单位向量 \(\bs{a}\)，\(\bs{b}\) 的夹角为 \(60^{\circ}\)，\(\bs{c} = t\bs{a} + (1 - t)\bs{b}\)，若 \(\bs{b} \cdot \bs{c} = 0\)，则 \(t =\)\fillinblank{}.
\end{problem}

\begin{answer}
2
\end{answer}

\begin{solution}
因为 \(\bs{a}\)，\(\bs{b}\) 是单位向量且夹角为 \(60^\circ\)，所以 \(\bs{b}\cdot\bs{a}=\cos60^\circ=\frac{1}{2}\)，\(\bs{b}\cdot\bs{b}=1\). 于是
\(\bs{b}\cdot\bs{c}=t\bs{b}\cdot\bs{a}+(1-t)\bs{b}\cdot\bs{b}=\frac{t}{2}+1-t=1-\frac{t}{2}\). 由 \(\bs{b}\cdot\bs{c}=0\) 得 \(t=2\).
\end{solution}

\begin{problem}
设 \(x\)，\(y\) 满足约束条件 \(\begin{cases}
1 \leqslant x \leqslant 3,\\
-1 \leqslant x - y \leqslant 0,
\end{cases}\)，则 \(z = 2x - y\) 的最大值为\fillinblank{}.
\end{problem}

\begin{answer}
3
\end{answer}

\begin{solution}
由 \(-1\leqslant x-y\leqslant0\) 得 \(x\leqslant y\leqslant x+1\). 对固定的 \(x\)，\(z=2x-y\) 随 \(y\) 增大而减小，因此取最小的 \(y=x\) 时，\(z=x\). 再由 \(1\leqslant x\leqslant3\)，得 \(z\) 的最大值为 \(3\).
\end{solution}

\begin{problem}
已知 \(H\) 是球 \(O\) 的直径 \(AB\) 上一点，\(AH:HB = 1:2\)，\(AB \perp\) 平面 \(\alpha\)，\(H\) 为垂足，平面 \(\alpha\) 截球 \(O\) 所得截面的面积为 \(\pi\)，则球 \(O\) 的表面积为\fillinblank{}.
\end{problem}

\begin{answer}
\(\frac{9\pi}{2}\)
\end{answer}

\begin{solution}
设球半径为 \(R\). 因为 \(AH:HB=1:2\) 且 \(AB=2R\)，所以 \(AH=\frac{2R}{3}\). \(O\) 是 \(AB\) 的中点，故 \(OH=R-\frac{2R}{3}=\frac{R}{3}\). 截面面积为 \(\pi\)，截面半径为 \(1\)，由球的截面关系得
\(\displaystyle R^2-OH^2=1^2\)，即 \(R^2-\frac{R^2}{9}=1\)，所以 \(R^2=\frac{9}{8}\). 球的表面积为 \(4\pi R^2=\frac{9\pi}{2}\).
\end{solution}

\begin{problem}
设当 \(x = \theta\) 时，函数 \(f(x) = \sin x - 2\cos x\) 取得最大值，则 \(\cos \theta =\)\fillinblank{}.
\end{problem}

\begin{answer}
\(-\frac{2\sqrt{5}}{5}\)
\end{answer}

\begin{solution}
由柯西不等式，
\(\displaystyle f(x)=\sin x-2\cos x\leqslant\sqrt{1^2+(-2)^2}\sqrt{\sin^2x+\cos^2x}=\sqrt{5}\). 取得最大值时，\((\sin\theta,\cos\theta)\) 与 \((1,-2)\) 同方向，因此 \(\sin\theta=\frac{1}{\sqrt{5}}\)，\(\cos\theta=-\frac{2}{\sqrt{5}}=-\frac{2\sqrt{5}}{5}\).
\end{solution}

\section{解答题}

\begin{problem}
已知等差数列 \(\{a_{n}\}\) 的前 \(n\) 项和 \(S_{n}\) 满足 \(S_{3} = 0\)，\(S_{5} = -5\).
\begin{enumerate}
\item 求 \( \{a_{n}\} \) 的通项公式；
\item 求数列 \(\left\{\frac{1}{a_{2n-1}a_{2n+1}}\right\}\) 的前 \(n\) 项和.
\end{enumerate}
\end{problem}

\begin{solution}
\begin{enumerate}
\item 设等差数列首项为 \(a_1\)，公差为 \(d\). 由
\(S_3=\frac{3}{2}\left(2a_1+2d\right)=3(a_1+d)=0\)，得 \(a_1+d=0\). 又
\(S_5=\frac{5}{2}\left(2a_1+4d\right)=5(a_1+2d)=-5\)，得 \(a_1+2d=-1\). 两式相减可得 \(d=-1\)，从而 \(a_1=1\)，所以 \(a_n=a_1+(n-1)d=2-n\).
\item 由 \(a_{2k-1}=3-2k\)，\(a_{2k+1}=1-2k\)，数列的第 \(k\) 项为
\(\displaystyle\frac{1}{(3-2k)(1-2k)}=\frac{1}{(2k-3)(2k-1)}=\frac{1}{2}\left(\frac{1}{2k-3}-\frac{1}{2k-1}\right)\). 因而前 \(n\) 项和为
\(\displaystyle\sum_{k=1}^{n}\frac{1}{a_{2k-1}a_{2k+1}}=\frac{1}{2}\left(-1-\frac{1}{2n-1}\right)=\frac{n}{1-2n}\).
\end{enumerate}
\end{solution}

\begin{problem}
为了比较两种治疗失眠症的药（分别称为 A 药，B 药）的疗效，随机地选取 \(20\) 位患者服用 A 药，\(20\) 位患者服用 B 药，这 \(40\) 位患者在服用一段时间后，记录他们日平均增加的睡眠时间（单位：\(\mathrm{h}\)）. 试验的观测结果如下：

服用 A 药的 \(20\) 位患者日平均增加的睡眠时间：

\(0.6\) \(1.2\) \(2.7\) \(1.5\) \(2.8\) \(1.8\) \(2.2\) \(2.3\) \(3.2\) \(3.5\)
\(2.5\) \(2.6\) \(1.2\) \(2.7\) \(1.5\) \(2.9\) \(3.0\) \(3.1\) \(2.3\) \(2.4\)

服用 B 药的 \(20\) 位患者日平均增加的睡眠时间：

\(3.2\) \(1.7\) \(1.9\) \(0.8\) \(0.9\) \(2.4\) \(1.2\) \(2.6\) \(1.3\) \(1.4\)
\(1.6\) \(0.5\) \(1.8\) \(0.6\) \(2.1\) \(1.1\) \(2.5\) \(1.2\) \(2.7\) \(0.5\)

\begin{enumerate}
\item 分别计算两组数据的平均数，从计算结果看，哪种药的疗效更好？
\item 根据两组数据完成下面茎叶图，从茎叶图看，哪种药的疗效更好？
\end{enumerate}

\begin{center}
\begin{tabular}{c|c|c}
A药 &  & B药 \\
\hline
\(6\) & \(0\) & \(5\) \(5\) \(6\) \(8\) \(9\) \\
\(8\) \(5\) \(5\) \(2\) \(2\) & \(1\) & \(1\) \(2\) \(2\) \(3\) \(4\) \(6\) \(7\) \(8\) \(9\) \\
\(9\) \(8\) \(7\) \(7\) \(6\) \(5\) \(4\) \(3\) \(3\) \(2\) & \(2\) & \(1\) \(4\) \(5\) \(6\) \(7\) \\
\(5\) \(2\) \(1\) \(0\) & \(3\) & \(2\)
\end{tabular}
\end{center}
\end{problem}

\begin{solution}
\begin{enumerate}
\item A 药数据的总和为 \(46.0\)，B 药数据的总和为 \(32.0\)，所以两组数据的平均数分别为
\(\displaystyle\overline{x}_A=\frac{46.0}{20}=2.3\)，\(\displaystyle\overline{x}_B=\frac{32.0}{20}=1.6\). 从平均数看，A 药疗效更好.
\item 以整数部分为茎、十分位数字为叶，叶按从小到大的顺序排列. A 药各茎的叶分别为：茎 \(0\) 的叶为 \(6\)，茎 \(1\) 的叶为 \(2\)，\(2\)，\(5\)，\(5\)，\(8\)，茎 \(2\) 的叶为 \(2\)，\(3\)，\(3\)，\(4\)，\(5\)，\(6\)，\(7\)，\(7\)，\(8\)，\(9\)，茎 \(3\) 的叶为 \(0\)，\(1\)，\(2\)，\(5\). B 药各茎的叶分别为：茎 \(0\) 的叶为 \(5\)，\(5\)，\(6\)，\(8\)，\(9\)，茎 \(1\) 的叶为 \(1\)，\(2\)，\(2\)，\(3\)，\(4\)，\(6\)，\(7\)，\(8\)，\(9\)，茎 \(2\) 的叶为 \(1\)，\(4\)，\(5\)，\(6\)，\(7\)，茎 \(3\) 的叶为 \(2\). 这些叶与题中给出的 \(20\) 个数据逐一对应. A 药有更多数据落在茎 \(2\)、\(3\) 中，整体分布位置高于 B 药，因此从茎叶图看 A 药疗效更好.
\end{enumerate}
\end{solution}

\begin{problem}
如图，三棱柱 \(ABC - A_{1}B_{1}C_{1}\) 中，\(CA = CB\)，\(AB = AA_{1}\)，\(\angle BAA_{1} = 60^{\circ}\).

\begin{enumerate}
\item 证明：\(AB \perp A_{1}C\)
\item 若 \(AB = CB = 2\)，\(A_{1}C = \sqrt{6}\)，求三棱柱 \(ABC - A_{1}B_{1}C_{1}\) 的体积.
\end{enumerate}

\bitmapfigure[max width=0.46\linewidth]{img/2013/new_curriculum_paper_1_liberal/q19_fig1.png}
\end{problem}

\begin{solution}
\begin{enumerate}
\item 设 \(\overrightarrow{AB}=\bs{u}\)，\(\overrightarrow{AC}=\bs{v}\)，\(\overrightarrow{AA_1}=\bs{w}\). 由 \(CA=CB\) 得
\(\lvert\bs{v}\rvert=\lvert\bs{v}-\bs{u}\rvert\)，所以 \(\bs{u}\cdot\bs{v}=\frac{1}{2}\lvert\bs{u}\rvert^2\). 又 \(AB=AA_1\)，且 \(\angle BAA_1=60^\circ\)，所以
\(\bs{u}\cdot\bs{w}=\lvert\bs{u}\rvert\lvert\bs{w}\rvert\cos60^\circ=\frac{1}{2}\lvert\bs{u}\rvert^2\). 因此
\(\bs{u}\cdot\overrightarrow{A_1C}=\bs{u}\cdot(\bs{v}-\bs{w})=0\)，故 \(AB\perp A_1C\).
\item 取 \(A\) 为坐标原点，令 \(B=(2,0,0)\)，\(C=(1,\sqrt{3},0)\). 因为 \(AA_1=AB=2\) 且 \(\angle BAA_1=60^\circ\)，可设 \(A_1=(1,u,v)\)，并由 \(AA_1=2\) 得 \(u^2+v^2=3\). 又 \(A_1C=\sqrt{6}\)，所以
\(\displaystyle6=(1-1)^2+(\sqrt{3}-u)^2+v^2=3-2\sqrt{3}u+u^2+v^2=6-2\sqrt{3}u\)，从而 \(u=0\)，再得 \(\lvert v\rvert=\sqrt{3}\). 于是三棱柱的高为 \(\sqrt{3}\). 底面 \(ABC\) 为边长 \(2\) 的等边三角形，面积为 \(\frac{\sqrt{3}}{4}\times2^2=\sqrt{3}\)，故体积为 \(\sqrt{3}\times\sqrt{3}=3\).
\end{enumerate}
\end{solution}

\begin{problem}
已知函数 \(f(x) = \e^{x}(ax + b) - x^{2} - 4x\)，曲线 \(y = f(x)\) 在点 \((0,f(0))\) 处的切线方程为 \(y = 4x + 4\).
\begin{enumerate}
\item 求 \(a\)，\(b\) 的值；
\item 讨论 \( f(x) \) 的单调性，并求 \( f(x) \) 的极大值.
\end{enumerate}
\end{problem}

\begin{solution}
\begin{enumerate}
\item 由切线经过 \((0,f(0))\)，得 \(f(0)=4\)，即 \(b=4\). 又
\(\displaystyle f'(x)=\e^x(ax+a+b)-2x-4\)，所以 \(f'(0)=a+b-4\). 切线斜率为 \(4\)，故 \(a+b-4=4\). 代入 \(b=4\) 得 \(a=4\).
\item 代入 \(a=b=4\)，得
\(\displaystyle f'(x)=\e^x(4x+8)-2x-4=2(x+2)(2\e^x-1)\). 因此临界点为 \(x=-2\) 和 \(x=-\ln2\). 当 \(x<-2\) 时，\(x+2<0\) 且 \(2\e^x-1<0\)，故 \(f'(x)>0\)；当 \(-2<x<-\ln2\) 时，\(f'(x)<0\)；当 \(x>-\ln2\) 时，\(f'(x)>0\). 所以 \(f(x)\) 在 \((-\infty,-2)\) 和 \((-\ln2,+\infty)\) 上单调递增，在 \((-2,-\ln2)\) 上单调递减. 在 \(x=-2\) 处导数由正变负，故函数取得极大值
\(\displaystyle f(-2)=\e^{-2}(-8+4)-4-4(-2)=4(1-\e^{-2})\).
\end{enumerate}
\end{solution}

\begin{problem}
已知圆 \(M: (x + 1)^2 + y^2 = 1\)，圆 \(N: (x - 1)^2 + y^2 = 9\)，动圆 \(P\) 与圆 \(M\) 外切并与圆 \(N\) 内切，圆心 \(P\) 的轨迹为曲线 \(C\).

\begin{enumerate}
\item 求 \(C\) 的方程；
\item \(l\) 是与圆 \(P\)，圆 \(M\) 都相切的一条直线，\(l\) 与曲线 \(C\) 交于 \(A\)，\(B\) 两点，当圆 \(P\) 的半径最长时，求 \(|AB|\).
\end{enumerate}
\end{problem}

\begin{solution}
\begin{enumerate}
\item 设圆 \(P\) 的半径为 \(r\)，圆心坐标为 \(P(x,y)\). 圆 \(P\) 与圆 \(M\) 外切、与圆 \(N\) 内切，故
\(\displaystyle PM=r+1\)，\(PN=3-r\). 两式相加得 \(PM+PN=4\). 圆 \(M(-1,0)\) 与圆 \(N(1,0)\) 是椭圆的两个焦点，因此 \(P\) 的轨迹满足 \(\displaystyle PM+PN=4\)，是以原点为中心、长半轴为 \(2\)、焦半距为 \(1\) 的椭圆. 于是短半轴平方为 \(4-1=3\)，轨迹方程为
\(\displaystyle\frac{x^2}{4}+\frac{y^2}{3}=1\). 当 \(P=(-2,0)\) 时 \(r=PM-1=0\)，不构成圆，故还应排除 \(x=-2\).
\item 由椭圆上 \(PM\) 的最大值为 \(3\)，可知圆 \(P\) 半径最长时 \(r=2\)，此时 \(P=(2,0)\). 对于这两个圆，公共切线有两类. 第一类是两圆外切点处的公共切线 \(x=0\)，它与椭圆交于 \((0,\sqrt3)\) 和 \((0,-\sqrt3)\)，所以 \(|AB|=2\sqrt3\).

另一类公共外切线可设为 \(x+ky+4=0\). 由其到 \(M(-1,0)\) 和 \(P(2,0)\) 的距离分别为 \(1\) 和 \(2\)，得 \(k^2=8\)，故可取 \(x\pm2\sqrt2y+4=0\). 以 \(x+2\sqrt2y+4=0\) 为例，代入椭圆方程得
\(\displaystyle\frac{(-2\sqrt2y-4)^2}{4}+\frac{y^2}{3}=1\)，即 \(7y^2+12\sqrt2y+9=0\). 两交点的纵坐标之差为 \(\frac{6}{7}\)，由直线方程可知横坐标之差为 \(2\sqrt2\cdot\frac{6}{7}\)，所以弦长为
\(\displaystyle\frac{6}{7}\sqrt{1+\left(2\sqrt2\right)^2}=\frac{6}{7}\sqrt9=\frac{18}{7}\). 因而 \(|AB|=2\sqrt3\) 或 \(\frac{18}{7}\).
\end{enumerate}
\end{solution}

\begin{problem}
如图，直线 \(AB\) 为圆的切线，切点为 \(B\)，点 \(C\) 在圆上，\(\angle ABC\) 的角平分线 \(BE\) 交圆于点 \(E\)，\(DB\) 垂直 \(BE\) 交圆于 \(D\).

\begin{enumerate}
\item 证明： \(DB = DC\)；
\item 设圆的半径为 \(1\)，\(BC = \sqrt{3}\)，延长 \(CE\) 交 \(AB\) 于点 \(F\)，求 \(\triangle BCF\) 外接圆的半径.
\end{enumerate}

\bitmapfigure[max width=0.42\linewidth]{img/2013/new_curriculum_paper_1_liberal/q22_fig1.png}
\end{problem}

\begin{solution}
\begin{enumerate}
\item 设 \(\angle ABE=\angle EBC=\alpha\). 由切线与弦所夹的角等于弦所对的圆周角，\(\angle BCE=\angle ABE=\alpha\)，所以 \(\triangle BCE\) 为等腰三角形，且 \(CE=BE\). 又 \(DB\perp BE\)，故 \(\angle DBE=90^\circ\)，从而 \(DE\) 是圆的直径. 在圆周上，\(\angle EDB=\angle ECB=\alpha\)，所以在直角三角形 \(BDE\) 中，\(\angle BED=90^\circ-\alpha\). 同弦 \(BD\) 所对的圆周角相等，故 \(\angle BCD=\angle BED=90^\circ-\alpha\). 另一方面，\(\angle DBC=90^\circ-\angle EBC=90^\circ-\alpha\). 因此 \(\angle DBC=\angle BCD\)，从而 \(DB=DC\).
\item 圆的半径为 \(1\)，弦 \(BC=\sqrt3\)，所以弦 \(BC\) 所对的圆心角为 \(120^\circ\). 由切线与弦所夹的角，\(\angle ABC=60^\circ\)，从而角平分线给出 \(\angle ABE=\angle EBC=30^\circ\). 又由切线与弦所夹的角，\(\angle BCE=\angle ABE=30^\circ\). 因为 \(F\) 在 \(CE\) 的延长线上且 \(F\) 在切线 \(AB\) 上，所以 \(\angle CBF=\angle CBA=60^\circ\)，\(\angle BCF=\angle BCE=30^\circ\)，于是 \(\angle BFC=90^\circ\). 三角形 \(BCF\) 为直角三角形，斜边 \(BC=\sqrt3\)，其外接圆直径为 \(BC\)，故外接圆半径为 \(\frac{\sqrt3}{2}\).
\end{enumerate}
\end{solution}

\begin{problem}
已知曲线 \(C_1\) 的参数方程为 \(\begin{cases}
x = 4 + 5\cos t,\\
y = 5 + 5\sin t,
\end{cases}\)（\(t\) 为参数），以坐标原点为极点，\(x\) 轴的正半轴为极轴建立极坐标系，曲线 \(C_2\) 的极坐标方程为 \(\rho = 2\sin\theta\).
\begin{enumerate}
\item 把 \(C_1\) 的参数方程化为极坐标方程；
\item 求 \(C_{1}\) 与 \(C_{2}\) 交点的极坐标 \(\left(\rho \geqslant 0, 0 \leqslant \theta < 2\pi\right)\).
\end{enumerate}
\end{problem}

\begin{solution}
\begin{enumerate}
\item 由参数方程消去 \(t\)，得 \((x-4)^2+(y-5)^2=25\)，即 \(x^2+y^2-8x-10y+16=0\). 代入 \(x=\rho\cos\theta\)，\(y=\rho\sin\theta\)，得到
\(\rho^2-8\rho\cos\theta-10\rho\sin\theta+16=0\).
\item 方程 \(\rho=2\sin\theta\) 等价于 \(x^2+y^2=2y\). 与 \(C_1\) 的直角坐标方程联立，作差得 \(x+y=2\). 代入 \(x^2+y^2=2y\)，得
\(\displaystyle(2-y)^2+y^2=2y\)，即 \(y^2-3y+2=0\). 因此交点为 \((1,1)\) 和 \((0,2)\). 在 \(\rho\geqslant0\)，\(0\leqslant\theta<2\pi\) 的条件下，它们的极坐标分别为 \(\left(\sqrt2,\frac{\pi}{4}\right)\) 和 \(\left(2,\frac{\pi}{2}\right)\).
\end{enumerate}
\end{solution}

\begin{problem}
已知函数 \( f(x) = |2x - 1| + |2x + a| \)，\( g(x) = x + 3 \).
\begin{enumerate}
\item 当 \(a = -2\) 时，求不等式 \(f(x) < g(x)\) 的解集；
\item 设 \(a > -1\)，且当 \(x \in \left[-\frac{a}{2}, \frac{1}{2}\right)\) 时，\(f(x) \leqslant g(x)\)，求 \(a\) 的取值范围.
\end{enumerate}
\end{problem}

\begin{solution}
\begin{enumerate}
\item 当 \(a=-2\) 时，\(f(x)=|2x-1|+|2x-2|\). 按 \(x=\frac{1}{2}\)，\(x=1\) 分段：
\[
f(x)=
\begin{cases}
3-4x, & x<\frac{1}{2},\\
1, & \frac{1}{2}\leqslant x<1,\\
4x-3, & x\geqslant1.
\end{cases}
\]
在 \(x<\frac{1}{2}\) 时，\(3-4x<x+3\) 等价于 \(x>0\)，得到 \(\left(0,\frac{1}{2}\right)\)；在 \(\frac{1}{2}\leqslant x<1\) 时，不等式恒成立；在 \(x\geqslant1\) 时，\(4x-3<x+3\) 等价于 \(x<2\)，得到 \([1,2)\). 合并得解集为 \((0,2)\).
\item 因为 \(a>-1\)，所以 \(-\frac{a}{2}<\frac{1}{2}\). 当 \(x\in\left[-\frac{a}{2},\frac{1}{2}\right)\) 时，\(2x+a\geqslant0\)，\(2x-1<0\)，从而 \(f(x)=1-2x+2x+a=a+1\). 要使 \(f(x)\leqslant g(x)=x+3\) 对该区间内所有 \(x\) 成立，只需在左端点取最小值，即
\(\displaystyle a+1\leqslant-\frac{a}{2}+3\). 解得 \(a\leqslant\frac{4}{3}\). 与 \(a>-1\) 合并，得 \(a\in\left(-1,\frac{4}{3}\right]\).
\end{enumerate}
\end{solution}
